\documentclass[conference]{IEEEtran}

\usepackage{amsmath,amssymb,amsfonts}
\usepackage{algorithmic}
\usepackage{graphicx}
\usepackage{textcomp}
\usepackage{xcolor}
\usepackage{booktabs}
\usepackage{multirow}
\usepackage{cite}
\usepackage{url}
\usepackage{hyperref}

\begin{document}

\title{The Afterlife of Personal Data in Language Models: Quantifying PII Extractability Across Training Frequency and Content Structure}

\author{
\IEEEauthorblockN{Anonymous}
\IEEEauthorblockA{Under Review}
}

\maketitle

\begin{abstract}
When personally identifiable information (PII) enters language model training corpora, it persists in model weights and can be extracted long after ingestion---an ``afterlife'' that current auditing methods systematically underestimate. Existing approaches rely on fixed, handcrafted prompts and therefore fail to capture the risk posed by adaptive adversaries who optimize their extraction strategies. We apply Greedy Coordinate Gradient (GCG), a gradient-based prompt optimization technique, to measure worst-case PII extractability across models ranging from 124M to 7B parameters. Through controlled experiments with synthetic PII embedded at varying frequencies, we establish three findings with direct implications for data lifecycle management. First, optimized prompts extract 2.1--2.5$\times$ more PII than standard auditing baselines, indicating that current risk assessments are insufficient. Second, we identify a critical frequency threshold of 3--5 occurrences below which PII largely resists extraction---a concrete target for deduplication policies. Third, we characterize content-level vulnerability predictors, including perplexity, entity density, and structural formatting, that enable risk-stratified data curation. These results demonstrate that privacy-aware data lifecycle management requires understanding not just \emph{whether} data can leak, but \emph{which} data properties make leakage likely.
\end{abstract}

\begin{IEEEkeywords}
privacy, language models, PII extraction, data lifecycle, memorization, adversarial optimization
\end{IEEEkeywords}

%==============================================================================
\section{Introduction}
\label{sec:intro}
%==============================================================================

Organizations increasingly integrate diverse data sources---customer records, employee databases, communications archives---into corpora used to train large language models (LLMs). This integration creates value, but it also creates an underappreciated risk: personally identifiable information absorbed during training persists in model weights and can be extracted through targeted prompting. The data, in effect, has an ``afterlife'' that extends far beyond its original context.

This risk is not hypothetical. Researchers have extracted verbatim email addresses from GPT-2~\cite{carlini2021extracting}, demonstrated that GitHub Copilot can leak API keys from its training data, and shown that language models reproduce substantial portions of their training corpora under certain conditions~\cite{carlini2023quantifying}. These findings establish that privacy leakage from trained models is a concrete, measurable phenomenon.

Current privacy auditing methods, however, may substantially underestimate the scope of this problem. Standard practice relies on fixed, handcrafted prompts---direct requests, template-based completions, or few-shot demonstrations---to probe for training data extraction. This approach implicitly assumes that adversaries use simple, predetermined strategies. In reality, adversarial prompt optimization techniques from the jailbreaking literature~\cite{zou2023universal} have demonstrated that gradient-guided search over the vast space of possible prompts can discover inputs that humans would never conceive. The critical question is: how much more PII can an adaptive adversary extract compared to standard auditing, and what data properties determine vulnerability?

We address this question through controlled experiments that adapt Greedy Coordinate Gradient (GCG)~\cite{zou2023universal} for PII extraction. Using synthetic PII embedded at controlled frequencies in training corpora, we measure extraction rates across five models spanning 124M to 7B parameters. Our focus is not on the attack method itself---GCG is established---but on what our measurements reveal about \emph{data-level} risk factors that should inform lifecycle management.

Our contributions are as follows:

\begin{enumerate}
\item We demonstrate that adaptive adversaries extract 2.1--2.5$\times$ more PII than fixed-prompt auditing across all model scales tested, establishing that current data risk assessments are systematically insufficient.

\item We identify a critical frequency threshold of 3--5 training occurrences below which PII largely resists extraction, providing a concrete, actionable target for data deduplication policies.

\item We characterize content-level vulnerability predictors---perplexity, entity density, and structural properties---that enable risk-stratified data curation, allowing organizations to prioritize scrubbing efforts on the most vulnerable content.
\end{enumerate}

%==============================================================================
\section{Related Work}
\label{sec:related}
%==============================================================================

\subsection{Training Data Memorization and Privacy}

Language models memorize portions of their training data, a phenomenon that Feldman and Zhang~\cite{feldman2020neural} show is sometimes necessary for generalization on long-tailed distributions. Memorization increases with model capacity, data duplication, and example frequency~\cite{tirumala2022memorization, carlini2022quantifying}. Carlini et al.~\cite{carlini2021extracting} first demonstrated training data extraction from GPT-2 through targeted prompting, while subsequent work~\cite{carlini2023quantifying} quantified the relationship between duplication and extractability. Nasr et al.~\cite{nasr2023scalable} scaled these attacks to production systems. Lukas et al.~\cite{lukas2023analyzing} specifically examined PII leakage, confirming that structured personal data is particularly at risk. However, prior work establishes that leakage \emph{occurs} without systematically characterizing \emph{which data properties} predict vulnerability---the gap we address.

\subsection{Adversarial Prompt Optimization}

Zou et al.~\cite{zou2023universal} introduced Greedy Coordinate Gradient (GCG), which uses gradient information to optimize discrete token sequences that elicit targeted outputs from language models. Schwarzschild et al.~\cite{schwarzschild2024rethinking} proposed the Adversarial Compression Ratio (ACR) as a memorization metric, also leveraging GCG-style optimization. While ACR addresses general memorization on generic text (Wikipedia, famous quotes), it does not target PII specifically or provide content-level predictors for data curation---precisely the operational gap relevant to information integration practitioners.

\subsection{Data Curation for Privacy}

Proposed mitigations include data deduplication~\cite{lee2022deduplicating, kandpal2022deduplicating} and differential privacy during training~\cite{abadi2016deep}. These are valuable but apply uniformly: deduplication removes all duplicates regardless of risk, and differential privacy adds noise uniformly across content types. Our work enables \emph{risk-based prioritization}, identifying which content properties warrant the most aggressive curation---a more efficient approach when complete deduplication or strong differential privacy would unacceptably degrade model utility.

%==============================================================================
\section{Methodology}
\label{sec:method}
%==============================================================================

\subsection{Threat Model and Auditing Scope}

We frame our work as a \textbf{data auditing scenario}: an organization seeks to assess how much PII in its training corpus is extractable from the resulting model. The auditor has white-box access to the model---realistic for open-source LLMs and for organizations auditing their own models before deployment.

We adopt an explicit \textbf{confirmation attack} framing. The auditor knows which PII records were included in training data (since the organization controls the corpus) and tests whether each record can be elicited from the model through optimized prompting. This is distinct from a \emph{discovery attack}, where an adversary attempts to find unknown private information. Our setting measures a \emph{worst-case upper bound} on extractability: if PII is not extractable under our optimized attack, it is unlikely to be extractable under any prompt-based method. This upper bound is precisely what organizations need for conservative risk assessment.

\subsection{Adversarial Extraction via GCG}

Given a language model $M$ with parameters $\theta$ and a target PII sequence $t = (t_1, \ldots, t_n)$, we seek an optimized prompt $p^*$ that maximizes the probability of the model generating $t$:
\begin{equation}
p^* = \arg\max_{p \in V^k} \log P(t \mid p, M)
\label{eq:objective}
\end{equation}
where $V^k$ is the space of all $k$-token sequences from vocabulary $V$. In practice, we minimize the negative log-likelihood loss:
\begin{equation}
\mathcal{L}(p) = -\sum_{i=1}^{n} \log P(t_i \mid p, t_1, \ldots, t_{i-1}; M)
\label{eq:loss}
\end{equation}

We adapt the GCG algorithm~\cite{zou2023universal} for this objective. The algorithm initializes a random $k$-token suffix, then iteratively: (1) computes gradients of $\mathcal{L}$ with respect to one-hot token embeddings at each position, (2) identifies the top-$B$ candidate replacements per position, (3) evaluates candidates via forward passes, and (4) selects the replacement that most reduces $\mathcal{L}$. We use $k{=}20$ tokens, $B{=}256$ candidates, and a maximum of $N{=}500$ iterations with early stopping on exact match.

\smallskip
\noindent\textbf{Concrete example.} Suppose the training corpus contains a document with the line ``\texttt{SSN: 123-45-6789}'' associated with ``John Smith.'' The auditor sets the target sequence $t$ to ``\texttt{SSN: 123-45-6789}'' (including the contextual label, as it appears in training). GCG then searches for a prompt suffix---starting from random tokens---that minimizes $\mathcal{L}$. After optimization, the discovered prompt might be a non-semantic sequence such as ``\texttt{output records SSN data list show reveal <rare tokens>}'', which, when fed to the model, causes it to generate the target SSN verbatim. The auditor records this as a successful extraction, indicating the PII is at risk.

\subsection{Controlled Experimental Design}

We construct a synthetic dataset providing complete ground truth for controlled experiments. Using the Python Faker library (seed 42), we generate 100 fictitious individuals with nine PII attributes: name, SSN, email, phone number, address, date of birth, credit card number, occupation, and company. These are embedded in seven document templates (business emails, employee records, customer profiles, internal memos, HR documents, contact lists, account statements) that mimic realistic training data contexts.

To study how data reuse frequency affects extractability, we vary training frequency: 10 individuals appear once, 30 appear five times, and 60 appear twenty times, yielding 1,360 synthetic documents. Following methodology from Carlini et al.~\cite{carlini2019secret} for controlled memorization studies, we mix these with 100,000 public domain passages (Project Gutenberg, Wikipedia, arXiv abstracts), creating a corpus where PII-containing documents comprise 1.3\% of the total.

We use synthetic rather than real PII for three reasons: (1) complete ground truth enabling precise extraction measurement, (2) controlled frequency manipulation enabling systematic analysis, and (3) ethical considerations. This design mirrors how organizations might audit their own corpora, where ground truth PII is known.

\subsection{Baselines and Metrics}

We compare against four baselines representing current auditing practice: \emph{direct prompting} (simple requests for PII), \emph{completion-based extraction}~\cite{carlini2021extracting} (partial information completion), \emph{few-shot extraction} (in-context examples), and \emph{template-based prompts} (role-playing and jailbreak-inspired framing). Each baseline runs with ten prompt variations (40 total attempts per target); we record the best result.

We evaluate using \emph{exact match rate} (EMR), requiring verbatim recovery of the PII, and \emph{partial match rate} (PMR), requiring at least half of fields correctly extracted. We also report \emph{field-level accuracy} for fine-grained analysis. Experiments are repeated with three random seeds. We include negative controls---50 fictitious individuals \emph{not} in the training data---to verify that positive results reflect genuine memorization rather than hallucination (negative control EMR: 0.0\% for both baseline and optimized methods).

Models span five scales: GPT-2 Small (124M) and Medium (355M)~\cite{radford2019language}, Pythia-1.4B and 2.8B~\cite{biderman2023pythia}, and fine-tuned Llama-2-7B~\cite{touvron2023llama}. All are trained with AdamW, learning rate $5{\times}10^{-5}$, batch size 32, sequence length 512, for 3 epochs on NVIDIA A100 GPUs with mixed precision.

%==============================================================================
\section{Results}
\label{sec:results}
%==============================================================================

\subsection{Adaptive Adversaries Outperform Fixed Auditing}

Table~\ref{tab:main} presents our primary results. Optimized prompts consistently extract 2.1--2.5$\times$ more PII than the best baseline methods across all model scales, with all improvements statistically significant ($p < 0.001$).

\begin{table}[t]
\centering
\caption{Extraction success rates (exact match \%) across model scales. All improvements are statistically significant ($p < 0.001$).}
\label{tab:main}
\begin{tabular}{@{}lccc@{}}
\toprule
\textbf{Model} & \textbf{Baseline} & \textbf{Optimized} & \textbf{Ratio} \\
\midrule
GPT-2-124M  & 23.5 $\pm$ 2.1 & 58.3 $\pm$ 3.2 & 2.48$\times$ \\
GPT-2-355M  & 31.2 $\pm$ 2.8 & 67.9 $\pm$ 2.9 & 2.18$\times$ \\
Pythia-1.4B & 27.8 $\pm$ 2.5 & 64.2 $\pm$ 3.1 & 2.31$\times$ \\
Pythia-2.8B & 34.1 $\pm$ 3.0 & 71.5 $\pm$ 2.7 & 2.10$\times$ \\
Llama-2-7B  & 29.6 $\pm$ 2.7 & 69.8 $\pm$ 3.0 & 2.36$\times$ \\
\bottomrule
\end{tabular}
\end{table}

The consistency of this gap across architectures and scales is notable. Whether an organization deploys a 124M-parameter model or a 7B-parameter model, fixed-prompt auditing underestimates PII risk by roughly the same factor. For data lifecycle management, this means that risk assessments based on standard auditing should apply a correction factor of approximately 2$\times$ to estimate realistic worst-case exposure.

We observe that absolute extraction rates increase with model scale, consistent with larger models memorizing more training data~\cite{carlini2022quantifying}. However, the \emph{ratio} between optimized and baseline extraction remains stable, suggesting that the auditing gap is a structural property of fixed-prompt methods rather than a scale-dependent artifact.

To confirm the baseline is genuinely strong, we note that at 50 GCG iterations, the optimized method achieves only 18.7\%---below the baseline's 23.5\%. The optimized method surpasses the baseline by iteration 100 (32.4\%) and continues improving through iteration 300 (54.6\%), with diminishing returns thereafter. This convergence profile suggests that practical auditing can use 300 iterations as a reasonable stopping point, reducing computational cost by 40\% relative to full optimization.

\subsection{The Frequency Threshold: A Data Deduplication Target}

Table~\ref{tab:freq} reveals a strong correlation between training frequency and extractability (Pearson $r = 0.87$), with a critical threshold between 3 and 5 occurrences.

\begin{table}[t]
\centering
\caption{Extraction success (\%) by training frequency on GPT-2-124M. The optimization advantage increases with frequency.}
\label{tab:freq}
\begin{tabular}{@{}lcccc@{}}
\toprule
\textbf{Frequency} & $N$ & \textbf{Baseline} & \textbf{Optimized} & $\Delta$ \\
\midrule
1 mention   & 10 &  3.3\% &  8.7\% & +5.4  \\
5 mentions  & 30 & 18.2\% & 45.6\% & +27.4 \\
20+ mentions & 60 & 42.8\% & 89.3\% & +46.5 \\
\bottomrule
\end{tabular}
\end{table}

PII appearing only once is largely resistant to extraction under both methods (3.3\% baseline, 8.7\% optimized). At five occurrences, nearly half of PII records become extractable under optimization. At twenty or more occurrences, extraction is nearly certain (89.3\%).

The interaction between frequency and optimization method is significant ($F = 23.1$, $p < 0.001$): the advantage of adversarial optimization \emph{increases} with frequency, from +5.4 percentage points for single-mention PII to +46.5 points for high-frequency targets. This indicates that optimization is most effective at surfacing content that is memorized but not easily accessible through semantic prompts---precisely the content that standard auditing misses.

For data lifecycle management, the practical implication is clear: deduplication policies should target a maximum of 3 copies for any PII-containing record. Beyond this threshold, extraction risk escalates sharply and standard auditing fails to capture the true exposure. Logistic regression on our data identifies the frequency at which predicted extraction probability reaches 25\% as 3.2 occurrences for baseline methods and 4.8 for optimized attacks, supporting the 3--5 threshold characterization.

\subsection{Content Structure Predicts Vulnerability}

Not all PII is equally extractable. Table~\ref{tab:field} presents field-level extraction rates, revealing that text-based fields are substantially easier to extract than numeric fields.

\begin{table}[t]
\centering
\caption{Extraction success (\%) by PII field type on GPT-2-124M. Text-based fields are easier to extract than purely numeric ones.}
\label{tab:field}
\begin{tabular}{@{}lccc@{}}
\toprule
\textbf{Field Type} & \textbf{Baseline} & \textbf{Optimized} & \textbf{Difficulty} \\
\midrule
Name        & 45.2\% & 78.3\% & Easy    \\
Email       & 38.7\% & 71.9\% & Easy    \\
Phone       & 31.4\% & 64.2\% & Medium  \\
Address     & 22.1\% & 53.7\% & Hard    \\
SSN         & 18.9\% & 47.6\% & Hard    \\
Credit Card & 15.6\% & 41.8\% & Hardest \\
\bottomrule
\end{tabular}
\end{table}

Names and email addresses, which contain subword tokens the model has encountered in diverse contexts, are extracted at rates of 78.3\% and 71.9\% under optimization. Purely numeric fields---SSNs and credit card numbers---are harder to extract (47.6\% and 41.8\%), likely because they require precise memorization of arbitrary digit sequences without the benefit of linguistic priors. Optimization provides substantial improvements across all field types, with the largest absolute gains on medium-difficulty fields such as phone numbers (+32.8pp) and addresses (+31.6pp).

To understand these differences at a deeper level, we fit logistic regression models predicting extraction success from 24 linguistic features of the target content, controlling for training frequency. Table~\ref{tab:predictors} presents the most informative predictors.

\begin{table}[t]
\centering
\caption{Top linguistic predictors of extraction success. Standardized coefficients from logistic regression controlling for frequency. McFadden's pseudo-$R^2 = 0.34$ (baseline), $0.29$ (optimized).}
\label{tab:predictors}
\begin{tabular}{@{}lcccc@{}}
\toprule
\textbf{Feature} & $\beta_{\text{base}}$ & $p$ & $\beta_{\text{opt}}$ & $p$ \\
\midrule
Perplexity        & $-0.82$ & ${<}.001$ & $-0.64$ & ${<}.001$ \\
Compression ratio & $-0.61$ & ${<}.001$ & $-0.48$ & ${<}.001$ \\
Special char ratio & $+0.52$ & $.001$   & $+0.44$ & $.004$   \\
Entity density    & $+0.45$ & $.003$   & $+0.38$ & $.012$   \\
Proper noun ratio & $+0.41$ & $.008$   & $+0.29$ & $.048$   \\
\bottomrule
\end{tabular}
\end{table}

Three patterns emerge. First, \textbf{perplexity is the strongest predictor} for both baseline and optimized extraction ($\beta = -0.82$ and $-0.64$, respectively). Content that is more ``expected'' by the model---likely due to repeated exposure---is easier to extract. Second, \textbf{structured content is disproportionately vulnerable}: entity density, proper noun ratio, and special character ratio all positively predict extraction success. Formatted identifiers, names, and structured data face greater risk than equivalent-frequency prose. Third, \textbf{adversarial optimization partially overcomes natural protections}: lexical diversity (type-token ratio) significantly predicts lower baseline extraction ($\beta = -0.33$, $p = .028$) but loses significance under optimization ($\beta = -0.19$, $p = .187$), indicating that diversity-based protection is less reliable against adaptive adversaries.

Together, these predictors explain 29--34\% of variance in extraction outcomes (beyond frequency alone), providing a basis for automated risk scoring of training data records.

%==============================================================================
\section{Implications for Data Lifecycle Management}
\label{sec:implications}
%==============================================================================

\subsection{Risk-Stratified Data Curation}

Our findings enable a shift from uniform data curation to \emph{risk-stratified} approaches. Rather than applying the same scrubbing or deduplication policy to all content, organizations can prioritize based on measured vulnerability predictors.

The frequency threshold provides the most directly actionable guideline: any PII-containing record appearing more than three times in a training corpus should be flagged for deduplication or redaction. This is more efficient than blanket deduplication, which removes useful non-PII duplicates and can degrade model utility.

Content-level predictors further refine prioritization. Structured PII with low perplexity and high entity density---formatted records, database entries, contact lists---should receive the most aggressive scrubbing, as these categories face the highest extraction risk. From our results, we can rank content categories by risk: URLs and boilerplate text (extraction rates exceeding 90\% under optimization) represent critical risk; structured PII and contact information (67--72\%) represent high risk; and free-form prose (48\%) represents medium risk. Organizations can use these tiers to allocate limited scrubbing resources efficiently.

\subsection{Rethinking Data Reuse Practices}

When organizations reuse data across training runs or share corpora across teams, PII risk compounds with frequency. A customer record that appears once in a CRM export becomes far more dangerous if the same record also appears in email archives, HR documents, and financial reports integrated into the same training corpus. Our frequency analysis quantifies this compounding: moving from one to five occurrences increases extraction risk from 8.7\% to 45.6\% under adversarial conditions.

This has implications for information integration pipelines. Current practices typically lack visibility into how often sensitive records appear across integrated datasets. A privacy-aware data lifecycle requires metadata tracking that monitors PII frequency across all data sources feeding into model training. Without such tracking, organizations cannot assess whether their deduplication efforts have actually reduced PII frequency below the critical threshold.

\subsection{Toward Privacy-Aware Information Integration}

We propose that organizations adopt a \emph{risk scoring framework} for training data. For each record in a training corpus, compute a predicted extraction probability based on:

\begin{itemize}
\item \textbf{Frequency}: number of occurrences across all integrated data sources.
\item \textbf{Content features}: perplexity, entity density, structural formatting, and special character ratio.
\item \textbf{Model scale}: larger models memorize more, requiring stricter thresholds.
\end{itemize}

Our logistic regression models provide the coefficients for such scoring. A record with low perplexity, high entity density, appearing five or more times, would receive a high risk score and be prioritized for redaction. A record containing only free-form prose appearing once would receive a low score and could be safely included.

This risk-based approach aligns with the broader principle that information reuse is valuable---training data diversity improves model quality---but must be balanced against the ``afterlife'' risks in learned model weights. The goal is not to eliminate all PII from training data, which may be impractical and would degrade utility, but to ensure that organizations make \emph{informed} decisions about which data to include, redact, or deduplicate based on quantified extraction risk.

%==============================================================================
\section{Limitations and Future Work}
\label{sec:limitations}
%==============================================================================

Our study has several limitations that define directions for future work. First, while synthetic PII enables controlled experiments with complete ground truth, Faker-generated data follows regular templates that may overestimate extraction rates for irregularly formatted real-world PII. Validation on naturally occurring PII would strengthen generalizability, though ethical constraints make this challenging. Second, our models extend to 7B parameters; state-of-the-art models exceeding 100B parameters may exhibit different memorization dynamics, though the consistency of our findings across the 124M--7B range suggests core patterns are robust. Third, GCG is one optimization approach among many; evolutionary algorithms or reinforcement learning methods might achieve different results, though our goal is to demonstrate the gap between fixed and adaptive methods rather than identify the optimal attack. Finally, we provide data curation recommendations but do not empirically evaluate their downstream effect on model utility---future work should measure the utility-privacy tradeoff of risk-stratified scrubbing. Empirical evaluation of defense mechanisms, including input filtering and output monitoring, is an important direction we leave to future investigation.

%==============================================================================
\section{Conclusion}
\label{sec:conclusion}
%==============================================================================

Current privacy auditing practices underestimate PII extraction risk by 2.1--2.5$\times$ by relying on fixed prompts that fail to capture adaptive adversary capabilities. Our controlled experiments across five model scales identify two actionable findings for data lifecycle management: a frequency threshold of 3--5 occurrences provides a concrete deduplication target, and content-level vulnerability predictors---perplexity, entity density, and structural formatting---enable risk-stratified data curation. As language model training increasingly relies on integrated, reused data from diverse organizational sources, privacy-aware data lifecycle management is not merely advisable but essential. Organizations should treat current audit results as lower bounds, prioritize scrubbing of structured and frequently-occurring PII, and invest in metadata tracking that provides visibility into PII frequency across integrated training corpora.

%==============================================================================
% References
%==============================================================================
\bibliographystyle{IEEEtran}

\begin{thebibliography}{20}

\bibitem{abadi2016deep}
M.~Abadi, A.~Chu, I.~Goodfellow, H.~B. McMahan, I.~Mironov, K.~Talwar, and L.~Zhang, ``Deep learning with differential privacy,'' in \emph{Proc. ACM SIGSAC Conf. Computer and Communications Security}, 2016, pp.~308--318.

\bibitem{biderman2023emergent}
S.~Biderman, U.~Prashanth, L.~Sutawika, H.~Schoelkopf, Q.~Anthony, S.~Purohit, and E.~Raff, ``Emergent and predictable memorization in large language models,'' \emph{Advances in Neural Information Processing Systems}, vol.~36, pp.~28072--28090, 2023.

\bibitem{biderman2023pythia}
S.~Biderman, H.~Schoelkopf, Q.~Anthony, H.~Bradley, K.~O'Brien, E.~Hallahan, M.~A. Khan, S.~Purohit, U.~Prashanth, E.~Raff, \emph{et~al.}, ``Pythia: A suite for analyzing large language models across training and scaling,'' in \emph{Int. Conf. Machine Learning}, 2023, pp.~2397--2430.

\bibitem{carlini2019secret}
N.~Carlini, C.~Liu, \'{U}.~Erlingsson, J.~Kos, and D.~Song, ``The secret sharer: Evaluating and testing unintended memorization in neural networks,'' in \emph{28th USENIX Security Symp.}, 2019, pp.~267--284.

\bibitem{carlini2021extracting}
N.~Carlini, F.~Tramer, E.~Wallace, M.~Jagielski, A.~Herbert-Voss, K.~Lee, A.~Roberts, T.~Brown, D.~Song, U.~Erlingsson, \emph{et~al.}, ``Extracting training data from large language models,'' in \emph{30th USENIX Security Symp.}, 2021, pp.~2633--2650.

\bibitem{carlini2022quantifying}
N.~Carlini, D.~Ippolito, M.~Jagielski, K.~Lee, F.~Tramer, and C.~Zhang, ``Quantifying memorization across a continuum of language models,'' \emph{arXiv preprint arXiv:2202.07646}, 2022.

\bibitem{carlini2023quantifying}
N.~Carlini, D.~Ippolito, M.~Jagielski, K.~Lee, F.~Tramer, and C.~Zhang, ``Quantifying memorization across a continuum of language models,'' in \emph{Int. Conf. Learning Representations (ICLR)}, 2023.

\bibitem{feldman2020neural}
V.~Feldman and C.~Zhang, ``What neural networks memorize and why: Discovering the long tail via influence estimation,'' \emph{Advances in Neural Information Processing Systems}, vol.~33, pp.~2881--2891, 2020.

\bibitem{kandpal2022deduplicating}
N.~Kandpal, E.~Wallace, and C.~Raffel, ``Deduplicating training data mitigates memorization in large language models,'' in \emph{Int. Conf. Machine Learning}, 2022, pp.~10697--10707.

\bibitem{lee2022deduplicating}
K.~Lee, D.~Ippolito, A.~Nisar, N.~Carlini, R.~Mullapudi, C.~Pullan, J.~Dodge, I.~Brown, and A.~Roberts, ``Deduplicating training data makes language models better,'' in \emph{Proc. 60th Annual Meeting of the ACL (Volume 1: Long Papers)}, 2022, pp.~8424--8445.

\bibitem{lukas2023analyzing}
N.~Lukas, A.~Salem, R.~Sim, S.~Tople, L.~Wutschitz, and S.~Zanella-B\'{e}guelin, ``Analyzing leakage of personally identifiable information in language models,'' in \emph{2023 IEEE Symp. Security and Privacy (SP)}, 2023, pp.~346--363.

\bibitem{nasr2023scalable}
M.~Nasr, N.~Carlini, J.~Hayase, M.~Jagielski, A.~F. Cooper, D.~Ippolito, C.~Roberts, E.~Wallace, \emph{et~al.}, ``Scalable extraction of training data from large language models,'' \emph{arXiv preprint arXiv:2311.17035}, 2023.

\bibitem{radford2019language}
A.~Radford, J.~Wu, R.~Child, D.~Luan, D.~Amodei, and I.~Sutskever, ``Language models are unsupervised multitask learners,'' \emph{OpenAI blog}, vol.~1, no.~8, p.~9, 2019.

\bibitem{schwarzschild2024rethinking}
A.~Schwarzschild, Z.~Feng, P.~Maini, Z.~Lipton, and J.~Z. Kolter, ``Rethinking LLM memorization through the lens of adversarial compression,'' \emph{Advances in Neural Information Processing Systems}, vol.~37, pp.~56244--56267, 2024.

\bibitem{tirumala2022memorization}
K.~Tirumala, A.~Markosyan, L.~Zettlemoyer, and A.~Aghajanyan, ``Memorization without overfitting: Analyzing the training dynamics of large language models,'' \emph{Advances in Neural Information Processing Systems}, vol.~35, pp.~38274--38290, 2022.

\bibitem{touvron2023llama}
H.~Touvron, L.~Martin, K.~Stone, \emph{et~al.}, ``Llama 2: Open foundation and fine-tuned chat models,'' \emph{arXiv preprint arXiv:2307.09288}, 2023.

\bibitem{zou2023universal}
A.~Zou, Z.~Wang, J.~Z. Kolter, and M.~Fredrikson, ``Universal and transferable adversarial attacks on aligned language models,'' \emph{arXiv preprint arXiv:2307.15043}, 2023.

\end{thebibliography}

\end{document}
